\documentclass[answers,12pt,addpoints]{exam}

\usepackage{import}
\import{C:/Users/prana/OneDrive/Desktop/MathNotes}{style.tex}

% Header
\newcommand{\name}{Pranav Tikkawar}
\newcommand{\course}{411 - Mathematical Analysis}
\newcommand{\assignment}{Homework 1}

\author{\name}
\title{\course\ - \assignment}

\begin{document}

\maketitle

\begin{questions}

\setcounter{question}{5}
\question
Prove that in a Hausdorff space (hence in every metric space),
every finite set is closed.

\begin{solution}
Let $X$ be Hausdorff and let $A\subseteq X$ be finite.
If $A=\varnothing$, then $A$ is closed.

Otherwise, write $A=\{a_1,\ldots,a_n\}$. We show that
$X\setminus A$ is open. Let $x\in X\setminus A$.
For each $i$, the Hausdorff property gives disjoint open sets
$U_i$ and $V_i$ with $x\in U_i$ and $a_i\in V_i$.

The set
\[
U=\bigcap_{i=1}^{n}U_i
\]
is open and contains $x$. Since $a_i\notin U_i$ for every $i$,
$U$ contains no point of $A$. Thus,
\[
x\in U\subseteq X\setminus A.
\]
Every point of $X\setminus A$ therefore has an open neighborhood
contained in $X\setminus A$. Hence, $X\setminus A$ is open,
so $A$ is closed.
\end{solution}


\setcounter{question}{7}
\question
Let $X$ be a metric space with metric $d$.

\begin{parts}

\part
Prove that the following function is also a metric on $X$:
\[
\overline{d}(p,q):=\min\{d(p,q),1\}.
\]

\begin{solution}
Let $p,q,r\in X$. We verify the metric axioms.

\begin{enumerate}[(i)]
\item Non-negativity:
Since $d(p,q)\geq 0$, we have $\overline{d}(p,q) = \min\{d(p,q),1\} \geq 0$.

\item Identity of indiscernibles:
Since $1>0$,
\[
\overline{d}(p,q)=0
\iff d(p,q)=0
\iff p=q.
\]

\item Symmetry:
Since $d$ is symmetric,
\[
\overline{d}(p,q)
=\min\{d(p,q),1\}
=\min\{d(q,p),1\}
=\overline{d}(q,p).
\]

\item Triangle inequality:
If $d(p,q)\geq 1$ or $d(q,r)\geq 1$, then
\[
\overline{d}(p,r)\leq 1
\leq \overline{d}(p,q)+\overline{d}(q,r).
\]
Otherwise, both distances are less than $1$, so
\[
\overline{d}(p,r)
\leq d(p,r)
\leq d(p,q)+d(q,r)
=\overline{d}(p,q)+\overline{d}(q,r).
\]
\end{enumerate}

Therefore, $\overline{d}$ is a metric on $X$.
\end{solution}


\part
Prove that $d$ and $\overline{d}$ induce the same topology;
that is, they have the same open sets.

\begin{solution}
For every $p\in X$ and $0<r\leq 1$,
\[
\overline{d}(p,q)<r
\iff \min\{d(p,q),1\}<r
\iff d(p,q)<r.
\]
Therefore,
\[
B_{\overline{d}}(p,r)=B_d(p,r).
\]

Suppose $U$ is open in $(X,d)$ and $p\in U$.
Choose $\epsilon>0$ such that $B_d(p,\epsilon)\subseteq U$,
and let $r=\min\{\epsilon,1\}$. Then
\[
B_{\overline{d}}(p,r)
=B_d(p,r)
\subseteq B_d(p,\epsilon)
\subseteq U.
\]
Thus, $U$ is open in $(X,\overline{d})$.

Conversely, suppose $U$ is open in $(X,\overline{d})$
and $p\in U$. Choose $\epsilon>0$ such that
$B_{\overline{d}}(p,\epsilon)\subseteq U$,
and let $r=\min\{\epsilon,1\}$. Then
\[
B_d(p,r)
=B_{\overline{d}}(p,r)
\subseteq B_{\overline{d}}(p,\epsilon)
\subseteq U.
\]
Thus, $U$ is open in $(X,d)$.

Hence, the two metrics induce the same topology.
\end{solution}

\end{parts}

\begin{remark}
Although the metrics induce the same topology, they may differ
in boundedness: $\overline{d}\leq 1$, even when $d$ is unbounded.
\end{remark}


\question
Let $X$ be an arbitrary topological space and let $A\subseteq X$.
Prove the following.

\begin{remark}[Partition Theorem]
Every point of $X$ lies in exactly one of the following:
the interior of $A$, the boundary of $A$, or the interior of
$X\setminus A$. Equivalently,
\[
X=\operatorname{int}(A)
\sqcup \partial A
\sqcup \operatorname{int}(X\setminus A).
\]
Here, $\sqcup$ denotes a disjoint union.
\end{remark}

\begin{parts}

\part
$A$ is open if and only if it contains none of its boundary points.

\begin{solution}
Suppose $A$ is open. For each $x\in A$, the set $A$ is an open
neighborhood of $x$ disjoint from $X\setminus A$. Thus,
$x\notin\partial A$, so $A\cap\partial A=\varnothing$.

Conversely, suppose $A\cap\partial A=\varnothing$.
If $x\in A$, then $x$ is neither a boundary point of $A$
nor an interior point of $X\setminus A$. By the Partition
Theorem, $x\in\operatorname{int}(A)$. Therefore, every point
of $A$ is an interior point, so $A$ is open.
\end{solution}


\part
$A$ is closed if and only if it contains all of its boundary points.

\begin{solution}
Suppose $A$ is closed. If $x\notin A$, then $X\setminus A$
is an open neighborhood of $x$ disjoint from $A$.
Thus, $x\notin\partial A$. Therefore,
$\partial A\subseteq A$.

Conversely, suppose $\partial A\subseteq A$.
If $x\in X\setminus A$, then $x$ is neither a boundary point
of $A$ nor an interior point of $A$. By the Partition
Theorem, $x\in\operatorname{int}(X\setminus A)$.
Thus, $X\setminus A$ is open, so $A$ is closed.
\end{solution}


\part
$A$ is neither open nor closed if and only if it contains
some of its boundary points, but not all of them.

\begin{solution}
By part (a), $A$ is not open if and only if
\[
A\cap\partial A\neq\varnothing.
\]
By part (b), $A$ is not closed if and only if
\[
\partial A\setminus A\neq\varnothing.
\]
Thus, $A$ is neither open nor closed exactly when it contains
at least one boundary point and omits at least one boundary point.
\end{solution}

\end{parts}


\question
Let $X$ be an arbitrary topological space. A set is called
\emph{clopen} if it is both closed and open.
Prove that $A\subseteq X$ is clopen if and only if
$\partial A=\varnothing$.

\begin{solution}
Suppose $A$ is clopen. By Problem 9(a), $A\cap\partial A=\varnothing$.
By Problem 9(b), $\partial A\subseteq A$. Together, these imply
$\partial A=\varnothing$.

Conversely, suppose $\partial A=\varnothing$. Then both
\[
A\cap\partial A=\varnothing
\qquad\text{and}\qquad
\partial A\subseteq A.
\]
By Problem 9(a) and (b), $A$ is both open and closed.
Hence, $A$ is clopen.
\end{solution}

\end{questions}

\end{document}